\documentclass[12pt,a4paper]{article}

\usepackage{fontspec}
\setmainfont{Libertinus Serif}
\usepackage[greek.ancient,spanish]{babel}
\babelfont[greek]{rm}{Libertinus Serif}
\usepackage[a4paper,margin=3.2cm]{geometry}
\usepackage{microtype}
\usepackage[hidelinks]{hyperref}

% Griego politónico escrito directamente en UTF-8:
% \gr{...} para palabras sueltas; bloques con \foreignlanguage{greek}{...}
\newcommand{\gr}[1]{\foreignlanguage{greek}{#1}}
\newcommand{\stanzabreak}{\par\vspace{0.6\baselineskip}}

\setlength{\parindent}{1.25em}
\setlength{\parskip}{0.35em}
\raggedbottom

\usepackage{xcolor}
\usepackage{amssymb}
\usepackage{enumitem}

\definecolor{rojocambio}{RGB}{170,25,25}
\newcommand{\quita}[1]{\textcolor{rojocambio}{\textit{[Suprimir: #1]}}}
\newcounter{cambio}

\newcommand{\propuesta}[5]{%
  \refstepcounter{cambio}%
  \subsection*{Cambio \thecambio: #1}
  \noindent\textbf{Localización:} #2.\par
  \noindent\textbf{Criterio:} #3\par
  \medskip
  \noindent\textbf{Original:}\par
  #4\par
  \medskip
  \noindent\textbf{Sugerencia:}\par
  #5\par
  \medskip
  \noindent$\square$ Aprobar\qquad$\square$ Rechazar\qquad$\square$ Revisar
  \par\bigskip
}

\hypersetup{
  pdftitle={Revisión de la integración de los poemas: Leer la Odisea en tiempos iletrados},
  pdfauthor={}
}

\title{Revisión de la integración de los poemas\\[0.4em]\large\emph{Leer la Odisea en tiempos iletrados}}
\author{}
\date{Agosto de 2026}

\begin{document}

\maketitle

\section*{Criterio}

Esta revisión comprueba si la prosa establece la función de cada poema sin explicarlo antes o repetirlo después. No propone suprimir ni reescribir versos. Tampoco considera redundantes las transiciones breves que cambian de escena, identifican al hablante o permiten seguir el desarrollo de un fragmento largo.

El resultado se limita a dos anticipaciones prescindibles. No se ha modificado el corpus.

\propuesta
  {No explicar dos veces la ausencia de introspección en Homero}
  {src/odisea11\_sirenas.tex, líneas 175--177}
  {La frase inmediatamente anterior ya afirma que Homero cuenta lo sucedido sin mostrar cómo se siente Odiseo ni en qué consiste la magia de las sirenas. El párrafo siguiente vuelve a formular la falta de introspección antes de que empiece el poema propio. Basta conservar la transición y anunciar la versión alternativa.}
  {«Esa era otra de las cuentas que yo tenía pendientes con el granbardo. En \emph{Abandonando Ítaca}, escribo mi propia versión. Homero no da muchos detalles de cómo se sentían los protagonistas de su historia. Los muestra a menudo comiendo o bebiendo ---generalmente en exceso---, discutiendo, llorando, vociferando, o intentando salvar el pellejo. Pero rara vez sentimos sus emociones, rara vez llegamos a ponernos en su piel. Y sin embargo, son hombres como nosotros».}
  {Esa era otra de las cuentas que yo tenía pendientes con el granbardo. En \emph{Abandonando Ítaca}, escribo mi propia versión\quita{. Homero no da muchos detalles de cómo se sentían los protagonistas de su historia. Los muestra a menudo comiendo o bebiendo ---generalmente en exceso---, discutiendo, llorando, vociferando, o intentando salvar el pellejo. Pero rara vez sentimos sus emociones, rara vez llegamos a ponernos en su piel. Y sin embargo, son hombres como nosotros}:}

\propuesta
  {Dejar que el poema amplíe el lamento de Circe}
  {src/odisea8\_circe.tex, línea 440}
  {El fragmento anterior muestra a Circe llorando por Elpénor. El siguiente ensancha el lamento hasta abarcar a todos los muertos de la guerra y del viaje. La prosa anuncia ya ese movimiento completo; si se limita a introducir la primera muerte, la ampliación queda en manos del poema.}
  {«Circe, en el recuerdo del viejo Ulises, deja de ser una diosa altiva, distante y no poco malvada, para convertirse en una mujer capaz de llorar por una muerte absurda, por todas las muertes absurdas que abundan en la \emph{Odisea}».}
  {Circe, en el recuerdo del viejo Ulises, deja de ser una diosa altiva, distante y no poco malvada, para convertirse en una mujer capaz de llorar por una muerte absurda\quita{, por todas las muertes absurdas que abundan en la \emph{Odisea}}:}

\section*{Elementos revisados y conservados}

No propongo intervenir en los siguientes casos:

\begin{itemize}[leftmargin=*]
  \item Los fragmentos de Nausícaa, Polifemo, Agamenón, los pretendientes y Penélope: la prosa sitúa el recuerdo o la culpa concreta y los versos desarrollan una lectura nueva.
  \item «La Gran Guerra»: no es una ilustración intercalada, sino un capítulo construido desde la conciencia del anciano Odiseo; las intervenciones en prosa orientan el avance entre sus distintas escenas.
  \item «El asesinato de las niñas»: el recorrido por Patroclo, Pentesilea, Héctor y Aquiles se aparta deliberadamente de la línea de la \emph{Odisea}, pero prepara el contraste moral entre Aquiles y Odiseo que cierra el capítulo.
  \item «Lecciones de conducir»: el poema autobiográfico permite explicar de manera concreta por qué el reencuentro entre Laertes y Odiseo deja frío al narrador; no funciona como adorno.
  \item «Marcharse de Ítaca»: los versos constituyen la resolución narrativa del libro. Las frases breves entre fragmentos marcan etapas distintas ---muerte, amor, regreso a Circe, liberación de los fantasmas y partida--- y forman una estructura deliberada.
  \item La presentación de la ampliación de \emph{Abandonando Ítaca} en «Menelao»: distingue el texto publicado del fragmento nuevo y proporciona al lector la situación narrativa antes de un pasaje extenso.
\end{itemize}

\end{document}
